\begin{table}[htbp]
\centering
\caption{Baseline latent-state model}
\label{tab:baseline_implied}
\begin{tabular}{lccc}
\toprule
 & (1) & (2) & (3)

 & No miscl. & Symmetric & Asymmetric
\midrule
Entry rate (\%) & 8.28 & 2.98 & 2.98

 & (0.05) & (0.06) & (0.06)

Exit rate (\%) & 8.54 & 2.91 & 2.91

 & (0.05) & (0.06) & (0.06)

False-positive rate (\%) & 0.00 & 2.99 & 3.05

 &  & (0.03) & (0.04)

False-negative rate (\%) & 0.00 & 2.99 & 2.92

 &  & (0.03) & (0.04)

Employment rate (\%) & 49.23 & 50.57 & 50.58

 & (0.18) & (0.53) & (0.53)

\multicolumn{4}{l}{\textit{Likelihood-ratio tests}}
$LR$: no misclassification & --- & 22,438.36 & ---

Boundary $p$-value & --- & $<0.001$ & ---

$LR$: symmetric errors & --- & --- & 7.49

$p$-value & --- & --- & 0.006

Log-likelihood & -357.204M & -349.371M & -349.368M

$N$ & 402,923 & 402,923 & 402,923

\bottomrule
\end{tabular}
\begin{minipage}{0.98\linewidth}
\footnotesize \textit{Note:} All specifications estimate the initial employment probability freely. Rates are percentages. Symmetric error imposes equal false-positive and false-negative probabilities; the asymmetric model reports them separately. Analytical SE in parentheses use an individual-level survey-weighted sandwich covariance matrix and the delta method. Employment is the model-implied steady-state share. LR statistics use weights normalized to sum to $N$. The test of no misclassification uses the 50:50 boundary mixture; the asymmetry test uses $\chi^2_1$. The calculation does not incorporate strata or PSU clustering.
\end{minipage}
\end{table}
